\documentclass[12pt]{article}
\usepackage[a4paper,margin=2.2cm]{geometry}
\usepackage{amsmath,amssymb,amsthm,mathtools}
\usepackage{hyperref}
\usepackage{cleveref}

% 中文（xelatex）
\usepackage{fontspec}
\usepackage{xeCJK}

\usepackage{fontspec}
\usepackage{xeCJK}
\setCJKmainfont{Noto Serif CJK SC}

\title{Linear Algebra}
\date{}
\begin{document}
\maketitle

\subsection*{Q1}
\begin{flushleft}
向量:\\
$\vec{a} = \begin{bmatrix}
    2 \\
    2 
\end{bmatrix},\vec{b} = \begin{bmatrix}
    3  \\
    0 
\end{bmatrix}$

1.计算 $\vec{a} \cdot \vec{b}$
\begin{align}
     \vec{a} \cdot \vec{b} = x1x2  +  y1y2 \\
      = 2\cdot 3 +  2 \cdot 0 
      = 6 
\end{align}

2.计算 $\cos \delta$
\begin{align}
     \vec{a} \cdot \vec{b} = |\vec{a}| |\vec{b}| \cos\delta \\
     = 6 = \sqrt{2^{2} + 2^{2}} \cdot \sqrt{3^{2} + 0^{2}} \cdot \cos\delta \\
     = 6 = \sqrt{72} \cdot \cos\delta \\
     = 6 = \sqrt{6^{2} \cdot 2} \cdot \cos\delta\\ 
     = 6 = 6\sqrt{2}\cdot \cos\delta\\ 
     cos\delta  = \frac{6}{6\sqrt{2}} \\
                = \frac{1}{\sqrt{2}} \\
                = \frac{\sqrt{2}}{2}
\end{align}

3.计算$\vec{a}$ 在 $\vec{b}$ 上的投影长度
\begin{align}
    = |\vec{a}| \cdot \frac{\sqrt{2}}{2}
    = 2\sqrt{2} \frac{\sqrt{2}}{2} \\ 
    = 2
\end{align}
\end{flushleft}


\subsection*{Q1}
\subsection*{写出逆时针旋转90度的矩阵R}
\begin{flushleft}
$R = \begin{bmatrix}
    1 & 0\\ 
    0 & 1
\end{bmatrix}$ \\ 
基向量 $i = [1 0], j = [0 1]$ \\
在标准R 中 i和j是标准90度夹角 \\ 
新的 基向量 $i = [0 1], j = [-1 0]$ \\
旋转90度之后 \\ 
$R = \begin{bmatrix}
    0 & -1 \\
    1 & 0 
\end{bmatrix}$
\end{flushleft}

\subsection*{Q3}
\subsection*{矩阵法解方程}
\begin{flushleft}
 方程组:
    \[
     \begin{cases}
        2x + y  = 3 \\
        x - y = 0 
     \end{cases}
     \]

1.写出A$\vec{x} = \vec{b}$
\[
\begin{bmatrix}
    2 & 1 \\
    1 & -1 
\end{bmatrix} @ \begin{bmatrix}
    x \\
    y
\end{bmatrix} = \begin{bmatrix}
    3 \\ 
    0
\end{bmatrix} 
\]
2.写出逆矩阵公式
\[
step1:
  2 \cdot -1 - 1 \cdot1 != 0  \text{ (exists an inverse matrix.)} 
\]
\[
  step2:
  \text{逆矩阵公式：}
\] 
$\frac{1}{ad - bc}\cdot \begin{bmatrix}
    d & -b \\
    -c & a
\end{bmatrix}$ \\ 

3.计算出x,y \\
$
= \frac{1}{-3} * \begin{bmatrix}
    -1 & -1 \\ 
    -1 & 2
\end{bmatrix} @ \begin{bmatrix}
    3 \\ 
    0
\end{bmatrix}
$\\
$
=  \frac{1}{-3} * \begin{bmatrix}
    -3 \\
    -3
\end{bmatrix}
$\\
$
= \begin{bmatrix}
    1 \\
    1
\end{bmatrix}
$
\end{flushleft}

\section*{Eigenvalue}
\subsection{Q1-Visual Inspection}
\begin{flushleft}
$
R = \begin{bmatrix}
    3 & 0 \\
    0 & 1
\end{bmatrix}
$ \\

1.向量$\vec{i}=\begin{bmatrix}
    1 & 0
\end{bmatrix}$ 变换后变成了什么？它还是x轴上吗 如果是 伸缩倍数$\lambda_{1}$是多少? \\
$
\vec{i}_{n} =
R @ \vec{i}
= 
\begin{bmatrix}
    3 \\
    0
\end{bmatrix}$ \\
$\vec{i}_{n} = \vec{i} \cdot \lambda_{1} ?$ \\ 
$\lambda_{1} = 3 $ \\

2.向量$\vec{j}=\begin{bmatrix}
    0 & 1
\end{bmatrix}$ 变换后变成了什么？它还是y轴上吗 如果是 伸缩倍数$\lambda_{2}$是多少? \\
$
\vec{j}_{n} =
R @ \vec{j}
= 
\begin{bmatrix}
    0 \\
    1
\end{bmatrix}$ \\ 
$\vec{j}_{n} = \vec{j} \cdot \lambda_{2} ?$ \\ 
$\lambda_{2} = 1 $ \\

1.向量$\vec{w}=\begin{bmatrix}
    1 & 1
\end{bmatrix}$ 它还在原来的方向（45度）上吗？（如果不在，说明它不是特征向量） \\
$
\vec{w}_{n} =
R @ \vec{w}
= 
\begin{bmatrix}
    3 \\
    1
\end{bmatrix}$ \\

确定斜率 = $\frac{3-0}{1-0}$  ($\vec{w}_{n}$ 和 ${0,0}$ 两个点) \\
$ = 3 \neq 1$ ($\vec{w}$的斜率)
证明 $\vec{w}_{n}$ 和 $\vec{w}$ 不在一个方向上,不是特征向量
\end{flushleft}


\subsection{Q2-Verification}
\begin{flushleft}
方向由什么决定？斜率 \\
是否存在$\lambda$ 的由斜率是否相同确定 \\
已知矩阵 $B = \begin{bmatrix}
    4 & 1 \\
    2 & 3
\end{bmatrix}$ \\ 
有人声称$\vec{v} = \begin{bmatrix}
    1 \\
    1
\end{bmatrix}$ 是它的特征向量。\\ 

1.请计算 $B\vec{v}$
$
= \begin{bmatrix}
    5 \\
    5
\end{bmatrix}
$ \\
2.与$\vec{v}的关系$ \\
1.step
斜率是否相同 5/5 == 1/1 ?
相等\\
5倍 \\

3.特征值$\lambda$ \\
$\lambda$ = 5
\end{flushleft}

\subsection*{Eigenvector}
\begin{flushleft}
基础公式 \\
$A\vec{u} = \lambda \vec{u}$ \\ 
移项 \\
$A\vec{v} - \lambda \vec{v} = 0$ \\ 
$(A - \lambda I)\vec{u} = 0$

\end{flushleft}

\subsection*{Eigenvector-1}
\begin{flushleft}
推导公式: \\
$A\vec{u} = \lambda\vec{u}$ \\
= 
$(A-\lambda I)\vec{u} =  0$ \\
= 
$M\vec{u} = 0$ \\




% 矩阵 A = $\begin{bmatrix}
%     4 & 1 \\
%     2 & 3
% \end{bmatrix}$ \\
% 1.写出矩阵 $A - \lambda I$. \\
% = $\begin{bmatrix}
%     4 & 1 \\
%     2 & 3
% \end{bmatrix} -  \begin{bmatrix}
%     \lambda & 0 \\
%     0 & \lambda  \\ 
% \end{bmatrix}$ \\ 
% = $\begin{bmatrix}
%    4 - \lambda & 1 \\
%    2  & 3 - \lambda \\
% \end{bmatrix}$ \\ 

% 2.计算它的行列式 $\det(A-\lambda I)$\\
% = $ ((4- \lambda)  * (3 - \lambda))-2 $ = 0 \\
% = $ (12 - 4\lambda - 3\lambda  + \lambda^{2}) - 2 = 0$\\
% = $ 10 - 7\lambda + \lambda^{2} = 0$\\
% = $ \lambda^{2} - 7\lambda + 10 = 0$\\
% = $ (\lambda-2) * (\lambda-5) = 0$\\
% = $\lambda = 2 \| \lambda = 5$ \\





% 我对特征向量和特征值的理解一直不是很清晰
% $A\vec{u} = \lambda\vec{u}$ \\
% 那么
% $A\vec{u} - \lambda\vec{u} = \vec{0}$


% 这里有个很有意思的地方 $M\vec{u}  = \vec{0}$
% 这里这个$\vec{0}$ as 0b 和上一个$\vec{0}$ as 0a 的意义是不同的
% 就如你所说的0a 是相似度 也就是$\lambda\vec{u}$ 和 $A\vec{u}$ 
% 在同一个地方?(或者我理解的有问题？) 而 0b 是把 $\vec{u}$ 扯回原点(0,0)??

% 所以这个M 是非常有意思的 首先它是一个矩阵 又基向量ij 组成 也就是说它对改变的环境
% 对所有向量都会生效 但如果存在一个向量$\vec{u}$被M 运算 变成了 $\vec{0}$ 则证明
% $\vec{u}$ 被"消失"了 而这个消失代表的什么意思呢 ？ 为什么这个消失会和行列式有关系
% 为什么又能表达特征性质呢？ 我可能有点笨 我无法理解


% 我继续问一些问题和使用我自己的语言表述下理解(我需要使用自己的语言简洁干净的表达出来 则证明我有可能理解了)
% $A\vec{u} = \lambda\vec{u}$ \\
% 上面这行公式其实隐藏了一个 $\begin{bmatrix}
%     1 & 0 \\
%     0 & 1  
% \end{bmatrix}$ \\ 
% = 
% $A\vec{u} - (\lambda \cdot \begin{bmatrix}
%     1 & 0 \\
%     0 & 1 
% \end{bmatrix})\vec{u} = \vec{0}$
% 这个$\vec{0}$ 来至两个向量的相减 它很清晰也就是代表着两个向量的位置是重叠的 \\ 
% (我这里使用的是重叠  也就是说 它不止是方向相同 而是位置完全相等 或者说我理解错了？) \\

% 我们把上面的公式变化下形态 \\
% = 
% $(A - \lambda \cdot \begin{bmatrix}
%     1 & 0 \\
%     0 & 1 
% \end{bmatrix})\vec{u} = \vec{0}$ \\ 
% = 
% $M\vec{u} = \vec{0}$ \\

% 上面的式子也很清晰 一个向量$\vec{u}$ 被矩阵M 运算变化之后的结果是 $\vec{0}$ 也就是原点(0) \\
% 那么我们会得到下面信息: \\
% 1: $\vec{u}$ 非0 (如果是0 就没有意义了) \\
% 2: $\vec{u}$ 完全被丢失了 找不回了 \\
% 3: M 的行列式为0 行列式为0的矩阵和向量运算时 才会存在"降维"甚至消失(压回到0点) \\
% 因为行列式的本质是 一个向量经过矩阵变化之后的体积变化率(这是二维的 三维是体积 依次) \\ 

% 那么我们仔细思考下M 到底是怎么诞生的 (A - \lambda \cdot \begin{bmatrix}
%     1 & 0 \\
%     0 & 1 
% \end{bmatrix}) 是两个矩阵相减的？
% 那么两个矩阵相减到底是什么几何性质呢 ？ + 呢 ？ 乘除呢 ？
\end{flushleft}


\begin{flushleft}
% 第一题（手算）：矩阵 A=[0−110]A = \begin{bmatrix} 0 & -1 \\ 1 & 0 \end{bmatrix}
% A=[01​−10​]，不写代码，用纸笔回答：


% e1⃗=[1,0]\vec{e_1} = [1, 0]
% e1​​=[1,0] 变换后去了哪里？

% e2⃗=[0,1]\vec{e_2} = [0, 1]
% e2​​=[0,1] 变换后去了哪里？

% 这个变换在几何上是什么操作？
% 向量 [3,2][3, 2]
% [3,2] 变换后是什么？（用"基向量去哪了"的思路算，不要用矩阵乘法公式硬算）

$\vec{e1} = [1, 0] -> [0, 1]$ \\
$[1,0]$ 在二维是标准的右边运动 \\
$[0,1]$ 是正上方运动 \\

$\vec{e2} = [0, 1] -> [-1, 0] $
$[0,1]$ 是正上方运动 \\
$[-1,0]$ 在二维是标准的左边运动 \\

这个变换在几何上是什么操作？
所以这个变化就是把向量 逆时针旋转到第二象限 90度

向量
[3,2] as v 变换后是什么？（用"基向量去哪了"的思路算，不要用矩阵乘法公式硬算）
$A\vec{v} = 
3 \cdot \vec{e1} +  2 \cdot \vec{e2}$ \\
\begin{flushleft}
    已知 $\vec{v1} = [a1, b1, c1] 和 \vec{v2} = [a2, b2, c2]$，\\
    求一个$\vec{n} = [x, y, z] 使用n 同时垂直于v1和v2$ \\
    
    "同时垂直于两个", 就是两个点积都等于 0: \\ 
    $$
    \left\{ 
    \begin{aligned}
        a1x + b1y + c1z = 0 \\ 
        a2x + b2y + c2z = 0
    \end{aligned} 
    \right.$$

\end{flushleft}

= $ \begin{bmatrix}
    0 \\
    3
\end{bmatrix} + \begin{bmatrix}
    -2\\ 
     0
\end{bmatrix}$ \\
= $\begin{bmatrix}
    -2\\
    3
\end{bmatrix}$
\end{flushleft}


\begin{flushleft}
    已知 $\vec{v1} = [a1, b1, c1] 和 \vec{v2} = [a2, b2, c2]$，\\
    求一个$\vec{n} = [x, y, z] 使用n 同时垂直于v1和v2$ \\
    
    "同时垂直于两个", 就是两个点积都等于 0: \\ 
    $$
    \left\{ 
    \begin{aligned}
        a1x + b1y + c1z = 0 \\ 
        a2x + b2y + c2z = 0
    \end{aligned} 
    \right.$$

\end{flushleft}










\end{document}
